%%%%%%%%%%  scan idx 205  |  folio 192  %%%%%%%%%%
affines. Soit $E$ un $\mathcal O_S$-Module acyclique, et considérons la résolution $C^\bullet(E)$ de $E$ définie par le recouvrement $(U_i)$ (T.F. II 5.2). Deux cas sont à examiner :

1°) $S$ est \emph{séparé}. Alors les composantes de $C^\bullet(E)$ sont des sommes de termes du type $f_*f^*(E)$, où $f:U\longrightarrow S$ est l'inclusion d'un ouvert affine (donc un morphisme affine, $S$ étant séparé). On a
\[
 \mathbb R(Qf_*)f^*E\simeq \mathbb RQ\,\mathbb Rf_*f^*E\simeq \mathbb RQ\,f_*f^*E,
\]
puisque $f^*E$ est acyclique. On a d'autre part, d'après (3.4),
\[
 \mathbb R(Qf_*)f^*E\simeq \mathbb R(f_*Q)f^*E\simeq f_*\mathbb RQf^*E,
\]
puisque $f$ est affine. On en déduit
\[
 R^nQ(f_*f^*E)\simeq f_*R^nQ(f^*E)
\quad\text{pour tout }n\geq0,
\]
d'où d'après le cas affine,
\[
 R^nQ(f_*f^*E)=0\quad\text{pour }n>N.
\]
Comme $C^\bullet(E)$ est de longueur $N'$, on a finalement
\[
 R^nQ(E)=0\quad\text{pour }n>N+N',
\]
et l'on a gagné.

2°) $S$ est \emph{noethérien}. Les composantes de $C^\bullet(E)$ sont des sommes de termes du type $f_*f^*(E)$, où $f:U\longrightarrow S$ est l'inclusion d'un ouvert quasi-compact et \emph{séparé}. Comme dans le cas 1°), on a
